\section{Bernard's Degrees of Love and the Everlasting Sabbath}

When Cardinal Aimeric, Chancellor of the Roman Church, wrote to St.\ Bernard of Clairvaux asking why and how God ought to be loved, the Cistercian abbot answered in the opening paragraph of \emph{On Loving God} (\emph{De diligendo Deo}) with stark directness:

\begin{quote}
You wish me to explain for what reason and in what measure we should love God. I should say that God Himself is the motive of our love to Him, and that the measure of due love is to be without measure.\endnote{St.\ Bernard of Clairvaux, \emph{On Loving God} (\emph{De diligendo Deo}), ch.\ I (translated as \emph{Saint Bernard on the Love of God} by Marianne Caroline and Coventry Patmore, second edition, London: Burns and Oates, 1884, p.\ 3).}
\end{quote}

Why should love of God know no measure? Bernard unfolds a double title that binds creation and redemption into one debt of wonder. First, God has loved us gratuitously when we were unworthy and even his enemies, though he has no need of our goods (Ps 15:2 DR; Ps 16:2 Heb.; 1 John 4:10, 19); what greater gift could God bestow than himself, delivering up his only-begotten Son for our sake (\emph{On Loving God}, ch.\ I)? Second, even apart from the Cross, every human being owes God breath, light, bread, free will, reason, and the very capacity to seek virtue (ch.\ II). Citing St.\ Paul's question---``what hast thou that thou hast not received'' (1 Cor 4:7)---Bernard identifies the twin truths of creaturely wisdom: to know both \emph{what} we are in our dignity, and \emph{that we are not such of ourselves}.\endnote{Bernard of Clairvaux, \emph{On Loving God}, ch.\ II (Patmore 1884, pp.\ 6--8).}

In Chapter VII, Bernard confronts the suspicion that loving God for the sake of beatitude is merely a refined bargain. His distinction cuts straight to the bone: ``Love is a going forth of the soul, not a contract,'' and ``True love seeks no reward, but it merits one.'' Nobody pays a starving person to eat or hires a tender mother to nurse her infant; love finds its recompense not in an external wage, but in the beloved himself. By contrast, those who try to quiet the soul with finite possessions resemble a starving traveler opening his mouth wide to drink gusts of wind: earthly things puff the soul without feeding it. Only God answers the hunger he arouses: ``He prevents, He sustains, He fills thee; He kindles desire in thy heart, and He Himself is the Object of all desire.''\endnote{Bernard of Clairvaux, \emph{On Loving God}, ch.\ VII (Patmore 1884, pp.\ 29--30, 34--35).}

Because God is infinite and the creature is finite, our love can never equal his in magnitude. Yet the creature's return is \emph{whole} even while remaining \emph{unequal}: we cannot love as much as God loves, yet we can love with all that we are---whole heart, whole soul, whole strength---responding to infinite gift with unreserved self-gift.

\subsection*{The Four Degrees of Love}

Bernard knows human nature too honestly to imagine that anyone leaps from raw self-interest into pure disinterested charity overnight. In Chapters VIII through XI of \emph{On Loving God}, he maps the patient pedagogy by which grace leads a creature up four degrees of love.

\textbf{First: loving oneself for one's own sake} (ch.\ VIII). Weak, fleshly nature begins where it stands, caring first for its own survival and comfort. Grace does not despise this natural root; it disciplines it. When self-love threatens to overflow its banks into greed and voluptuousness, God erects the dyke of the second commandment: ``Thou shalt love thy neighbour as thyself.'' Sharing what one denies one's own disordered appetites widens the heart: ``Selfishness becomes benevolence by taking a wider range.'' Yet because we cannot love our neighbor rightly except in God, and because trials and frailty soon teach us that we can neither sustain ourselves nor do lasting good alone, tribulation drives the soul to cry out to its Maker.

\textbf{Second: loving God for one's own sake} (ch.\ IX). At this stage, the believer turns to God genuinely, yet primarily as rescuer, protector, and giver of necessities. Bernard does not sneer at this second degree: ``It is already something to feel the limits of his own capacity, to know what he cannot do without the help of God, and to keep right with Him who sustains his life and strength.'' Still, a person who loves God only for aid has not yet learned to love God for who he is.

\textbf{Third: loving God for God's own sake} (ch.\ IX). Repeated recourse to God in necessity slowly changes the heart that knocks at his door. Tasting his fidelity, mercy, and beauty, the soul discovers a sweetness far surpassing the favors it originally came asking for. Like the Samaritans who told the woman at the well that they now believed because they had heard Christ themselves, the believer says to self-interest:

\begin{quote}
It is not because of thee that we love the Lord, but we have tasted ourselves, and found how sweet He is. ... Herein we see such love as that which said, ``Give praise to the Lord, for He is good.'' He that loves God because He is good, not to him, but in Himself, loves God indeed for God, and not as he does of whom the Scripture saith: ``He will praise Thee when Thou shalt do well to him.''\endnote{Bernard of Clairvaux, \emph{On Loving God}, ch.\ IX (Patmore 1884, pp.\ 41--42).}
\end{quote}

Here obedience ceases to feel like a foreign tax; loving God truly, the soul readily loves all that belongs to him, keeping the precept of brotherly love willingly and without calculation.

\textbf{Fourth: loving even oneself only for God's sake} (chs.\ X--XI). Bernard calls this summit a high mountain where God is pleased to dwell, scarcely glimpsed for a fleeting moment in this mortal life before bodily weakness, daily cares, and the call of neighborly duty pull the soul back down to earth. In Chapter X, trying to name what this complete conformity of wills looks like, Bernard reaches for three classic patristic images---and immediately draws the doctrinal boundary that keeps Christian union forever distinct from pantheism:

\begin{quote}
To attain to this, is for the soul to be deified; as a small drop of water appears lost if mixed with wine, taking its taste and colour; and as, when plunged into a furnace, a bar of iron seems to lose its nature and assume that of fire; or as the air filled with the sun's beams seems rather to become light than to be illuminated. So it is with the natural life of the Saints; they seem to melt and pass away into the will of God. For if anything merely human remained in man, how then should God be all in all? It is not that human nature will be destroyed, but that it will attain another beauty, a higher power and glory.\endnote{Bernard of Clairvaux, \emph{On Loving God}, ch.\ X (Patmore 1884, p.\ 45).}
\end{quote}

Remarkably, Bernard pushes his fidelity to creaturehood one step further in Chapter XI. Even after death, when the separated souls of the saints enter the depths of eternal light, they have not yet reached the final completeness of the fourth degree---because the human soul naturally loves and longs for its own body. Tracing the Bridegroom's invitation in the Song of Songs, Bernard maps three courses of a single banquet of charity: ``In this life, eat; after death you shall drink; in the kingdom of the resurrection you shall be inebriated.'' In mortal flesh, the faithful eat the bread of obedient work; after death, separated from the body, the saints drink wine still mingled with milk---their natural tenderness for the body yet to rise; only at the general resurrection, when the glorified, incorruptible body is reunited to the soul, will nothing divide the heart any longer. Then the blessed drink pure new wine at the marriage feast of the Lamb, knowing ``satiety which is never satisfied; ardour which is insatiable, yet most calm and peaceful; the eternal and incomparable desire to possess, which arises from no want.''\endnote{Bernard of Clairvaux, \emph{On Loving God}, ch.\ XI (Patmore 1884, pp.\ 47--52). Bernard's insistence that resurrection completes the saint's joy harmonizes with Benedict XII's later definition in \emph{Benedictus Deus} (1336): the separated soul already enjoys the true, immediate beatific vision of the divine essence, while the resurrection of the body adds the consummation of the whole human person in glory.}

\subsection*{Deification Without Absorption}

Bernard's sentence---``It is not that human nature will be destroyed, but that it will attain another beauty, a higher power and glory''---guards the decisive threshold of Catholic eschatology. Scripture dares to call the redeemed ``partakers of the divine nature'' (2 Pet 1:4) and Aquinas dares to call the blessed \emph{deiform} (\emph{ST} I, q.\,12, a.\,5), because grace and glory truly communicate a created share in God's own life of knowing and loving. Yet union with God is always \emph{created participation}, never ontological fusion, absorption into an impersonal sea, identity of essence, or the erasure of personal distinctness.

The First Vatican Council's \emph{Dei Filius} (ch.\ 1 and Canons I.3--4) draws that boundary at the very threshold of creation: God is ``really and essentially distinct from the world,'' and finite creatures never share one substance or essence with God or emanate from his substance.\endnote{First Vatican Council, \emph{Dei Filius}, ch.\ 1 and Canons I.3--4 (Manning 1912, pp.\ 219, 232--233; Latin text: ``re et essentia a mundo distinctus'').} St.\ Augustine inscribes the same boundary in the closing chapter of \emph{The City of God} (XXII.30): ``For it is one thing to be God, another thing to be a partaker of God.''\endnote{St.\ Augustine, \emph{The City of God} XXII.30 (translated by Marcus Dods, Edinburgh: T.\ \& T.\ Clark, 1871, Project Gutenberg eBook 45305, vol.\ 2; Latin CSEL 40.2, ed.\ E.\ Hoffmann, 1900).} Even St.\ John's revelation that ``God is charity'' (1 John 4:8, 16) moves strictly in one direction: it tells us who the living God is in his Trinitarian fullness, and can never be reversed into ``love is God,'' as if human emotion, mutual sympathy, or romantic fervor were divine.

Far from erasing personhood, union with God perfects it. As Augustine explains in \emph{The City of God} XXII.30, the first free will given to Adam was the ability \emph{not to sin}; the final free will given in glory is the gift of \emph{not being able to sin}---not because freedom has been extinguished, any more than God lacks freedom because he cannot do evil, but because the will beholding infinite Goodness clings to him indefeasibly with undivided joy. Nor do the redeemed lose their personal history in amnesia. Like a skilled physician who understands diseases professionally without suffering them in his own body, the saints keep an intellectual remembrance of their past ills while every bitter sting of sensible experience is blotted out---``for if they were not to know that they had been miserable, how could they, as the Psalmist says, for ever sing the mercies of God?''

\subsection*{The Everlasting Sabbath}

At the end of twenty-two books surveying the fall of empires, the tragedy of pride, the patience of pilgrimage, and the hope of the resurrection, St.\ Augustine brings \emph{The City of God} (XXII.30) to rest on the biblical figure of the Sabbath. God rested on the seventh day from all his works and blessed it; in the consummation of the ages, the redeemed themselves become that seventh day, replenished with his blessing, fulfilling the psalmist's word, ``Be still and see that I am God'' (Ps 45:11 DR; Ps 46:10 Heb.). Where the serpent once tempted Adam and Eve to seize divinity by deserting God---``you shall be as Gods'' (Gen 3:5)---God fulfills the creature's destiny by giving himself to be participated in, when ``God may be all in all'' (1 Cor 15:28).

And then Augustine links the seventh day of creation's rest to the eighth day of Christ's resurrection---the Lord's Day, the Sunday whose weekly rhythm has shaped the Church's worship through all her pilgrimage---and closes with the cadence that gathers every human \emph{why} into peace:

\begin{quote}
[S]uffice it to say that the seventh shall be our Sabbath, which shall be brought to a close, not by an evening, but by the Lord's day, as an eighth and eternal day, consecrated by the resurrection of Christ, and prefiguring the eternal repose not only of the spirit, but also of the body. There we shall rest and see, see and love, love and praise. This is what shall be in the end without end. For what other end do we propose to ourselves than to attain to the kingdom of which there is no end?\endnote{Augustine, \emph{The City of God} XXII.30 (Dods 1871, vol.\ 2; CSEL 40.2, p.\ 670: ``haec tamen septima erit sabbatum nostrum, cuius finis non erit vespera, sed dominicus dies velut octavus aeternus, qui Christi resurrectione sacratus est, aeternam non solum spiritus, verum etiam corporis requiem praefigurans. Ibi vacabimus et videbimus, videbimus et amabimus, amabimus et laudabimus. Ecce quod erit in fine sine fine. Nam quis alius noster est finis nisi pervenire ad regnum, cuius nullus est finis?'').}
\end{quote}

\emph{Ibi vacabimus et videbimus, videbimus et amabimus, amabimus et laudabimus.} Freed from anxious striving, the blessed behold God directly, love him as they see him, and overflow in eternal praise---``and they shall see his face: and his name shall be on their foreheads'' (Apoc 22:4).

\medskip

So the three movements of the creature's life before God come home to one another. First, we exist only because the Creator, needing nothing, freely holds us in being out of nothing at every instant. Second, justice answers that gift with truthful worship, lawful obedience, and right order toward God and neighbor---recognizing what is owed before presuming to claim what is gratuitous. Third, charity receives what justice could never purchase and philosophy could never force: the living God's own invitation into filial friendship, sustained by hope through prayer and dryness, fed at the Eucharistic table, tested in the relief of affliction and the love of our neighbor, and crowned in the Sabbath that has no evening.

At the end of every human \emph{why} stands the Triune God who loved us before we could seek him, who desired the creature before the creature ever knew to desire his face. ``My heart hath said to thee: My face hath sought thee: thy face, O Lord, will I still seek'' (Ps 26:8 DR). And when the veil falls, the creature who began as dust held in being discovers that its final perfection is simply to be itself, transfigured in light, before the Face of its Friend.
